%% maestro2tex -- generated table; do not edit by hand.
\begin{table}[htbp]
  \centering
  \caption{Large-signal drive, from the step response of \texttt{SlewOpenLoopWBiasTB} into \SI{10}{\pico\farad}. Slew rate is taken between the \SI{1.25}{\volt} and \SI{2.25}{\volt} output levels --- the upper half of the swing, which every corner traverses inside the \SI{20}{\micro\second} window and which excludes the current spike at each edge. The slewing currents are the mean load current over that same interval, i.e.\ the same measurement expressed the other way ($I=C\,\mathrm{d}V/\mathrm{d}t$ at \SI{10}{\pico\farad}), and the peak edge currents are the momentary spikes as the input steps. \emph{None of the four current rows is a DC drive rating.} The bench steps the differential input by \(\pm\SI{1.25}{\volt}\) open-loop, which takes the PMOS input pair outside its common-mode range, so these are large-signal slam currents into a purely capacitive load. For DC drive see Table~\ref{tab:ctrlamp-drivemap}.}
  \label{tab:ctrlamp-largesignal}
  \begin{tabular}{lrrrrrrr}
    \toprule
    Parameter & Unit & tt & ff & fs & sf & ss & Spread \\
    \midrule
    Slew rate, rising & \si{\volt\per\micro\second} & 1.403 & 1.315 & 1.379 & 1.407 & 1.473 & 0.158 \\
    Slew rate, falling & \si{\volt\per\micro\second} & 0.352 & 0.154 & 0.371 & 0.311 & 0.506 & 0.351 \\
    Slewing current, rising edge & \si{\micro\ampere} & 14.03 & 13.15 & 13.79 & 14.08 & 14.73 & 1.58 \\
    Slewing current, falling edge & \si{\micro\ampere} & 3.52 & 1.54 & 3.71 & 3.11 & 5.06 & 3.51 \\
    Peak edge current, sourcing & \si{\micro\ampere} & 24.38 & 18.09 & 24.52 & 23.33 & 26.27 & 8.18 \\
    Peak edge current, sinking & \si{\micro\ampere} & 12.58 & 10.81 & 12.40 & 12.51 & 13.34 & 2.53 \\
    \bottomrule
  \end{tabular}
\end{table}
